\newpage
\subsection{Widerstandsmessung III: Glühfaden}\label{sec:01-3}

\begin{highlightblock}[Temperatur eines Glühfadens]
Temperatur bei $U=\{40, 80, 120, 160, 200, 240\}$ Volt
\begin{enumerate}[label=\textbf{\alph*:}, ref=\alph*]
    \item R-Messung $R_\mathrm{kalt}$ des Glühfadens bei Umgebungstemperatur mit Ohmmeter\label{itm:P3-1}
    \item T-Messung der Raumtemperatur mit Thermometer.\label{itm:P3-2}
    \item U-I-Messung $R_\mathrm{heiss}$ des Glühfadens bei Betriebstemperatur mit den obrigen Spannungen.\label{itm:P3-3}
    \item Berechnung von $R(\theta_0)$ für $\theta_0=\SI{20}{\celsius}$.\label{itm:P3-4}
    \item Berechnung der Glühfadentemperatur bei bei den obrigen Spannungen aus $R_\mathrm{heiss}$.\label{itm:P3-5}
\end{enumerate}
\end{highlightblock}\medskip

\marginref{\ref{itm:P3-1},~\ref{itm:P3-2}}
Messung wie oben beschreiben ergibt $R_\mathrm{kalt} = \SI{30.9}{\ohm}$ und
$\theta_\mathrm{RT} = \SI{24.1}{\celsius}$.

\marginref{\ref{itm:P3-3}}

\begin{wrapfigure}[8]{r}{7cm}
    \vspace{-20pt}
    \centering
    \begin{circuitikz}[american inductors]
        \ctikzset{inductors/coils=6}
        \draw (0,0) coordinate (GND) to [sV] ++(0,2) coordinate (V1) -- ++(1,0) to[L] ++(0,-2) -- (GND);
        \draw (2,0) to[vL, name=L2] ++(0,2);
        \draw (L2.wiper) -- (L2.wiper |- V1)
            to[ameter, i=$I_A$] ++(2,0) coordinate (k1) 
            to[vmeter, v=$U_V$, *-*] (k1 |- GND) -- (2,0);
        \draw (k1) -- ++(1,0) to[lamp, l=$R_\mathrm{W}$, i>^=$I_R$] ++(0,-2) -- (k1 |- GND);
    \end{circuitikz}
    \caption{Aufbau der Messung}\label{fig:13-aufbau}
\end{wrapfigure}

Mit einem Variablen Trafo wurden die verschiedenen Spannungen eingestellt. Die
Schaltung wurde wie folgt aufgebaut: Die U-I-Messung wird wieder aufgrund des
niederohmigen Glühfadens mit einer Spannungsrichtigen Schaltung durchgeführt.
Der Widerstand ergibt sich somit wieder aus \autoref{eq:r-spannungsrichtig}.
Vor der Messung wird der erste Messbereich der Messgeräte für den
Kaltwiderstand wie in \autoref{sec:01-1} \autoref{itm:P1-2} abgeschätzt. Beim
ersten einschalten wird der höchste Strom erwartet, da der Glühdraht ein
Kaltleiter ist ($\alpha, \beta > 0$).

\begin{itemize}
    \item Voltmeter (M-4600): $U_V \sim \SI{40}{\volt}\implies$ MB: $\SI{200}{\volt}$
    \item Amperemeter (M-4630B): $I_A \sim \SI{1.3}{\ampere}\implies$ MB: $\SI{20}{\ampere}$
\end{itemize}

Sobald sich der Glühfaden erwärmt reduziert sich der Strom und der Messbereich
kann bei Bedarf angepasst werden. Die Messergebnisse (Effektivwerte) und der
daraus resultierende Widerstand sind in \autoref{tab:13-values} dargestellt.

\begin{table}[h]
    \centering
    \caption{Messergebnisse der Glühfadenmessung}\label{tab:13-values}
    \pgfplotstabletypeset[
        col sep=comma,
        header=true,
        every head row/.style={after row=\midrule},
        begin table={\begin{tabularx}{\textwidth}},
        end table={\end{tabularx}},
        columns/U/.style={column name={$U / \si{\volt}$}, column type=Y, precision=2, fixed},
        columns/UV/.style={column name={$U_V / \si{\volt}$}, column type=Y|, precision=2, fixed},
        columns/IA/.style={column name={$I_A / \si{\milli\ampere}$}, column type=Y, precision=3, fixed},
        columns/R/.style={column name={$R / \si{\ohm}$}, column type=Y, precision=2, fixed},
        columns/T/.style={column name={$T / \si{\celsius}$}, column type=Y, precision=2, fixed},
    ]{../data/13_Values.csv}
\end{table}

\marginref{\ref{itm:P3-4}}
Für die Temperaturabhängigkeit des Widerstands gilt die Näherung:

\begin{equation}
    R(\theta) = R(\theta_0) \cdot (1 + \alpha \cdot (\theta - \theta_0) + \beta \cdot (\theta - \theta_0)^2)
    \label{eq:gluehfaden}
\end{equation}

Für Wolfram sind die Temperaturkoeffizienten für $\theta_0=\SI{20}{\celsius}$:
$\alpha = \SI{4.1e-3}{\per\kelvin}$ und $\beta = \SI{1.0e-6}{\per\kelvin\squared}$.
Einsetzen der Messwerte aus \autoref{itm:P3-1} und \ref{itm:P3-2} in
\autoref{eq:gluehfaden} liefert:

\begin{equation*}
    R(\theta_0) = \frac{R_\mathrm{kalt}}{1 + \alpha \cdot (\theta_\mathrm{RT} - \theta_0) + \beta \cdot (\theta_\mathrm{RT} - \theta_0)^2} = \SI{30.39}{\ohm}
\end{equation*}

\marginref{\ref{itm:P3-5}}
Nun kann man die berechneten Widerstands Werte in die Linke Seite von \autoref{eq:gluehfaden}
einsetzen und mit einem Mathematik-Programm nach der Temperatur $\theta$ lösen. Die errechneten
Temperaturen wurden zu \autoref{tab:13-values} hinzugefügt.
